\paragraph{Two allocation estimands.}
For a specified sampling policy $\pi$ and learning procedure $f$, define
\begin{equation}
\Delta_{\pi}(T)=\mathbb{E}\!\left[S(f_{\pi,32,T})-S(f_{\pi,8,T})\right],
\label{eq:allocation_estimand}
\end{equation}
where $S$ is subject-balanced pooled macro-F1 and the expectation averages
training-subset randomization under the fixed cohort and evaluation protocol.
Natural sampling estimates a \emph{total allocation effect}, including label
distribution changes. The matched policy estimates a \emph{label-standardized
allocation effect} relative to the sampled $P=8$ quota vector, averaged over
these reference draws. It also imposes within-class donor round robin. These
are different policy targets, not competing estimators of a uniquely true
patient effect; neither isolates biological diversity from batch structure.

\paragraph{Stability is not uncertainty of the mean.}
The released joint stability interval is the central 95\% range obtained by
selecting one completed subset seed and bootstrapping paired test subjects.
It retains between-subset variability rather than averaging resampled seeds.
Its zero inclusion does not test $\Delta_{\pi}(T)=0$, and its positive-replicate
fraction is neither a posterior probability nor a significance test. Inference
for the seed-averaged effect requires a separately constructed interval;
the present main comparisons report stability, not that inferential claim.
